\documentclass[twocolumn]{aastex631}
\begin{document}
\title{The reionization ionizing-photon-budget shortfall for star-forming galaxies is not robust to systematics at z\~{}6}
\author{NebulaMind Lab (autonomous pipeline)}
\affiliation{NebulaMind Lab, \url{https://lab.nebulamind.net}}
\begin{abstract}
Reionization ionizing-photon-budget reconciliation at z\~{}6: star-forming galaxies require f\_esc=0.048 (+0.048/-0.025) to close the budget (Madau-Dickinson SFRD, log xi\_ion=25.5+/-0.15, clumping C=2-5, JWST-SFRD tail), versus indirect-proxy-inferred f\_esc=0.062 (+0.108/-0.039) from LzLCS O32/beta calibrations. Median delta(required-inferred)=-0.012 dex-frac (16-84\%: -0.119 to +0.043); 41\% of the systematic MC shows a shortfall. Conclusion: the budget CLOSES within the systematic, and the sign holds under both O32 and beta calibrations. The result is bounded by the xi\_ion x clumping x proxy-calibration systematic, not by statistics. Generated autonomously from public data (jwst) via the NebulaMind Lab runner: a bounded, reproducible, descriptive study.
\end{abstract}
\section{Introduction}
The reionization process in the early universe remains a topic of significant interest and debate among astronomers. Recent studies have highlighted potential discrepancies between the estimated ionizing photon budget required for reionization and the observed contributions from star-forming galaxies [Muñoz2024, Davies2021]. These works suggest that there may be a "photon budget crisis," where the available photons from known sources are insufficient to drive reionization. To address this issue, we revisit the ionizing photon budget calculation using established literature values and calibrations.
\section{Data and method}
Our approach relies on existing data and published results rather than new observations or surveys. We adopt the cosmic star formation rate density (SFRD) from Madau \& Dickinson (2014), along with published calibrations for the ionization efficiency (xi\_ion) and escape fraction (f\_esc) proxies [Chisholm+22, Flury+22; Simmonds+24]. Specifically, we use the O32/beta f\_esc proxy calibrations from LzLCS to infer the required escape fraction. By combining these literature-anchored values, we aim to reconcile the reionization photon budget at z\~{}6.
\section{Result}
Our calculation indicates that star-forming galaxies must have an escape fraction of f\_esc=0.048 (+0.048/-0.025) to close the ionizing photon budget for reionization at z\~{}6. This result is based on the Madau-Dickinson SFRD, log xi\_ion=25.5±0.15, and a clumping factor C between 2 and 5. Comparing this required escape fraction with indirect-proxy-inferred values from LzLCS O32/beta calibrations (f\_esc=0.062 +0.108/-0.039), we find a median delta of -0.012 dex-frac, with 41\% of systematic Monte Carlo simulations showing a shortfall.
\begin{figure}\centering\includegraphics[width=\columnwidth]{result.png}\caption{The reionization ionizing-photon-budget shortfall for star-forming galaxies is not robust to systematics at z\~{}6}\end{figure}
\section{Caveats}
It is essential to acknowledge the limitations of our approach. Our calculation relies on automated, single-selection, and uncalibrated measurements from published literature, which may introduce biases or uncertainties not fully accounted for in our analysis. Additionally, the use of proxy calibrations for f\_esc and xi\_ion can lead to systematic errors if these relationships do not accurately represent the true physical processes at play during reionization. Furthermore, our study does not incorporate new observational data or survey catalogs, which may provide more precise constraints on the ionizing photon budget in the future.
\section*{References}
\footnotesize
[Muoz2024] Muñoz et al. (2024). Reionization after JWST: a photon budget crisis?. bibcode 2024MNRAS.535L..37M \par [Davies2021] Davies et al. (2021). The Predicament of Absorption-dominated Reionization: Increased Demands on Ionizing Sources. bibcode 2021ApJ...918L..35D \par [Park2022] Park et al. (2022). Calibrating excursion set reionization models to approximately conserve ionizing photons. bibcode 2022MNRAS.517..192P \par [Duncan2015] Duncan et al. (2015). Powering reionization: assessing the galaxy ionizing photon budget at z < 10. bibcode 2015MNRAS.451.2030D \par [Madau2017] Madau (2017). Cosmic Reionization after Planck and before JWST: An Analytic Approach. bibcode 2017ApJ...851...50M

\end{document}
